\section{Idempotents}
In this section we determine idempotents $e_i$ using the following 2 facts:
\begin{enumerate}\itemsep=0pt
\item[1)] the metric, being invariant, is diagonal in the basis $\{e_1,e_2,e_3\}$,
\item[2)] the idempotents commute.
\end{enumerate}
Suppose that the following implicit cubic ODE def\/ines our $3$-web:
\begin{equation}\label{flat_cub}
p^3+S(x,y)p^2+A(x,y)p+B(x,y)=0.
\end{equation}
Let $p_1$, $p_2$, $p_3$ be the roots of~(\ref{flat_cub}) at a point
$(x,y)$ outside the discriminant curve. The following 1-forms vanish on the
solutions
\begin{gather}
\sigma_1=(p_2-p_3)(dy-p_1dx),\nonumber\\
\sigma_2=(p_3-p_1)(dy-p_2dx),\nonumber\\
\sigma_3=(p_1-p_2)(dy-p_3dx)\label{normalized forms}
\end{gather}
and satisfy
\[
\sigma_1+\sigma_2 + \sigma_3=0.
\]
Let us introduce an ``area'' form by
\[
\Omega =\sigma_1\wedge
\sigma_2=\sigma_2\wedge \sigma_3=\sigma_3\wedge
\sigma_1=(p_1-p_2)(p_2-p_3)(p_3-p_1)dy\wedge dx.
\]
 The Chern
connection form is def\/ined as (see~\cite{Be})
\[
\gamma:
=h_2\sigma_1-h_1\sigma_2=h_3\sigma_2-h_2\sigma_3=h_1\sigma_3-h_3\sigma_1,
\]
where $h_i$ verify the relations
\[
d\sigma_i=h_i\Omega.
\]
The web is f\/lat if\/f the connection form
is closed: $d(\gamma)=0$. This implies $d\sigma_i=\gamma \wedge \sigma_i$ and gives an integrating factor $k$ of the forms $\sigma_i$ as a solution to the following equation
\begin{gather*}%\label{dk}
dk=-\gamma k.
\end{gather*}
Computing the Chern connection form in terms of roots $p_i$ and using the Viete formulas one gets
\[
\gamma =\frac{\gamma_1dx+\gamma_2dy}{-D},
\]
where
\begin{gather*}
\gamma_1=\big(4BS^2-3AB-SA^2\big)S_x+B(SA-9B)S_y+(2A^2-6SB)A_x\nonumber\\
\hphantom{\gamma_1=}{}
+2B(3A-S^2)A_y+(9B-SA)B_x+\big(AS^2-4A^2+3SB\big)B_y,\nonumber\\
\gamma_2=\big(6SB-2A^2\big)S_x-2B\big(3A-S^2\big)S_y+(SA-9B)A_x\nonumber\\
\hphantom{\gamma_2=}{}
 + \big(4A^2-AS^2-3SB\big)A_y+\big(6A-2S^2\big)B_x+\big(2S^3+18B-8SA\big)B_y,%\label{connectionSAB}
\end{gather*}
and $D$ is the discriminant of~(\ref{flat_cub}) with respect to $p$
\[
D=18SAB+S^2A^2-4A^3-27B^2-4BS^3.
\]
We also need to know how the connection form in the normalization~(\ref{normalized forms})
is transformed when we change the local coordinates
\[
\bar{y}=f(x,y), \qquad \bar{x}=g(x,y).
\]
One f\/inds without dif\/f\/iculty that the forms $\bar{\sigma_i}$
in the new coordinates are related to the forms in old ones by
\begin{gather*}%\label{form transform}
\bar{\sigma_i}=\frac{(f_yg_x-f_xg_y)^2\sigma_i}{g_x^3-Sg_x^2g_y+Ag_xg_y^2-Bg_y^3}.
\end{gather*}
Therefore (see~\cite{Be})
\begin{equation}\label{connection transforms}
\bar{\gamma}=\gamma+d\ln\left(\frac{(f_yg_x-f_xg_y)^2}{g_x^3-Sg_x^2g_y+Ag_xg_y^2-Bg_y^3}\right).
\end{equation}
Let $v_i=\xi_i \partial_x+\eta_i\partial_y$ be commuting multi-valued vector f\/ields,
whose trajectories are the web leaves. They may have singularities at singular points:
with $K$ satisfying $\frac{dK}{K}=\gamma $ we have
\begin{gather*}%\label{v0}
v_1=\frac{K(\partial_x +p_1 \partial_y)}{(p_3-p_1)(p_1-p_2)},\qquad
v_2=\frac{K(\partial_x +p_2 \partial_y)}{(p_1-p_2)(p_2-p_3)},\qquad
v_3=\frac{K(\partial_x +p_3 \partial_y)}{(p_2-p_3)(p_3-p_1)}.
\end{gather*}
Here $x$, $y$ are restrictions on $S$ of f\/lat coordinates in $M$, in which the Euler vector f\/ield has the form $E=w_tt\partial_t+w_xx\partial_x+w_yy\partial_y$. A subtle point is that f\/lat coordinates are not necessarily the coordinates used for the normal forms: they are related to them by some coordinate transform preserving $X$.

The idempotents are
\begin{equation}\label{idempotents}
e_1=\alpha \partial_t+v_1, \qquad e_2=\beta \partial_t+v_2,\qquad e_1=(1-\alpha -\beta) \partial_t+v_3,
\end{equation}
where the components $\alpha$, $\beta$ are to be def\/ined by the orthogonality condition
\begin{equation}\label{orthogonal}
\langle e_1,e_2\rangle =\langle e_2,e_3\rangle =\langle e_3,e_1\rangle =0.
\end{equation}

$\bullet$ {\it Metric with $\langle e,e\rangle =0$.}
Equations~(\ref{orthogonal}) give
\begin{equation}\label{dfor0}
\delta=-1/K.
\end{equation} Adjusting the integration constant so that $K=-1$ we have $\delta =1$ and
\begin{equation}\label{alphabeta0}
\alpha =\frac{p_2p_3-p_1(p_2+p_3)}{2(p_1-p_2)(p_1-p_3)},\qquad \beta =\frac{p_1p_3-p_2(p_1+p_3)}{2(p_2-p_1)(p_2-p_3)}.
\end{equation}
Since $dK=0$, we see that $\gamma_1=\gamma_2=0$.
The equation $[e_1,e_2]=0$ reads as $v_1(\beta)=v_2(\alpha)$ and obviously implies $[e_1,e_3]=[e_2,e_3]=0$.
The above 3 equations $\gamma_1=\gamma_2=v_1(\beta)-v_2(\alpha)=0$
can be resolved with respect to $S_x$, $A_x$, $B_x$
to give the following system of hydrodynamic type
\begin{gather}\label{WDVV0}
S_x= \frac{1}{2}A_y,\qquad
A_x= 2By,\qquad
B_x= SB_y+BS_y-\frac{1}{2}AA_y.
\end{gather}
\begin{lemma}\label{web_ass_0}
Suppose that the solution web of ODE~\eqref{flat_cub} satisfies the following conditions:
\begin{enumerate}\itemsep=0pt
\item[$1)$] it admits an infinitesimal symmetry $X=w_xx\partial_x+w_yy\partial_y$ with $w_x\ne w_y$,
\item[$2)$] the coefficients $S$, $A$, $B$ satisfy the system of PDEs~\eqref{WDVV0}.
\end{enumerate}
Then there is a Frobenius $3$-fold germ with $\langle e,e\rangle =0$ whose booklet web is the solution web of~\eqref{flat_cub}.
\end{lemma}
\begin{proof}
Equations~(\ref{WDVV0}) ensures the local existence of a function $f(x,y)$ satisfying
$f_{yyy}=S$, $f_{yyx}=\frac{1}{2}A$, $f_{yxx}=B$, $f_{xxx}=BS-\frac{1}{4}A^2$.
Thus we obtain the associativity equation
\begin{equation}\label{assPoly}
f_{xxx}=f_{yyy}f_{yxx}-f^2_{yyx}.
\end{equation}
(See~\cite{MFa} where the equivalence of~(\ref{WDVV0}) and~(\ref{assPoly}) was established and~\cite{Al} where equations of these types were studied.)
Since the web admits an inf\/initesimal symmetry $X$, the function~$f$ is weighted homogeneous and def\/ines a germ of Frobenius $3$-fold with $\langle e,e\rangle =0$ (see~\cite{Df}). Equation~(\ref{flat_cub}) def\/ines the corresponding characteristic web of the above associativity equation. Hence the web is the booklet web of the constructed Frobenius $3$-fold germ (see~\cite{Aw}).
\end{proof}
The value of the lemma is that it makes unnecessary the checking the potentiality condition for the 4-tensor $(\nabla_z c)(u,v,w)$ in Def\/inition~\ref{defFrob}.
\begin{theorem}\label{web_Frobenius_0}
Suppose that the solution web of ODE~\eqref{flat_cub} satisfies the following conditions:
\begin{enumerate}\itemsep=0pt
\item[$1)$] it admits an infinitesimal symmetry $X=w_xx\partial_x+w_yy\partial_y$ with $w_x\ne w_y$,
\item[$2)$] its Chern connection form vanishes identically $\gamma=0$,
\item[$3)$] the vector fields $e_1$, $e_2$ commute, where $\alpha$, $\beta$ are as in~\eqref{alphabeta0}.
\end{enumerate}
Then there is a Frobenius $3$-fold germ with $\langle e,e\rangle=0$
whose booklet web is the solution web of~\eqref{flat_cub}.
\end{theorem}
\begin{proof}
Def\/ine a metric by~(\ref{metric0}) with $\delta=1$,
the multiplication as the direct sum of one-dimensional algebras with unities~(\ref{idempotents}),
and the Euler vector f\/ield by $E=w_tt\partial_t+X$ with $w_t=2w_y-w_x$.
Then Lemma~\ref{web_ass_0} ensures that the def\/ined structure is that of Frobenius \mbox{$3$-fold}.
\end{proof}
